\chapter{Architettura software}
\label{cap:software}

Questo capitolo descrive la libreria condivisa, il layer variazionale e le
componenti usate per appaiare le varianti; le procedure operative e i comandi
eseguiti sono raccolti nell'appendice~\ref{app:riproduzione}.

\section{Repository e libreria condivisa}
\label{sec:repository}

Il repository raccoglie sotto una radice comune la libreria condivisa
\code{quantum\_framework}, le quattro directory di esperimento, la suite di test e
la documentazione; la struttura completa è nell'appendice~\ref{app:riproduzione}.

Il file \file{run.py} seleziona l'esperimento e inoltra gli argomenti al relativo
\file{main.py}; la logica sperimentale rimane nelle quattro directory dedicate.

I quattro esperimenti, sviluppati inizialmente in repository separati, condividono
ora caricamento dei dati, addestramento, metriche e funzioni statistiche tramite
\code{quantum\_framework}. Le architetture e le varianti di ablazione restano
nelle rispettive directory. \classe{QuantumLinear}, usato negli esperimenti 01 e
03, appartiene alla libreria; l'esperimento 02 mantiene una propria
\classe{QuantumPatchEmbedding} con la stessa struttura circuitale.

I file \file{main.py} conservano la CLI e l'orchestrazione specifica.

Le dipendenze sono gestite da un unico \file{pyproject.toml} risolto con
\code{uv}, e l'ambiente pinnato in \file{uv.lock} è la fonte autoritativa. Il
requisito minimo è Python 3.12. Le librerie principali sono PennyLane per la
simulazione \parencite{bergholm2018pennylane}, PyTorch e torchvision per la parte
classica \parencite{paszke2019pytorch}, \code{medmnist} per i dataset,
\code{scikit-learn} e \code{scipy} per metriche e test statistici. Le run
definitive 01--03 usano il backend \code{default.qubit}; 01 e 02 forzano la CPU,
mentre nel 03 il corpo PyTorch gira su CUDA e il QNode resta sul simulatore.

\section{Il layer variazionale}
\label{sec:quantum-linear}

\classe{QuantumLinear} è usato dagli esperimenti 01 e 03 e presenta il circuito
come modulo PyTorch.
Non contiene alcun collo di bottiglia classico, per cui l'ultima dimensione
dell'ingresso deve coincidere con \code{n\_qubits}. Nell'esperimento 03 vale
$\dvit=\nq$; il modulo distinto dell'esperimento 02 applica lo stesso vincolo
(\cref{sec:base-config}). La pipeline, con le forme:
\begin{equation}
  \begin{aligned}
    \shape{\ldots, \nq}
    &\xrightarrow{\text{reshape}} \shape{M, \nq}
    \xrightarrow{\tanh \cdot \pi} [-\pi, \pi]^{\nq} \\
    &\to \texttt{AngleEmbedding}
    \to \texttt{StronglyEntanglingLayers}
    \to \big(\expval{Z_i}\big)_i
    \to \shape{\ldots, \nq}.
  \end{aligned}
  \label{eq:pipeline-vqc}
\end{equation}

La codifica angolare è periodica; prima del circuito gli input sono quindi
limitati con $\tanh$ e riscalati nell'intervallo $(-\pi,\pi)$. Per valori assoluti
elevati, $\tanh'(x)\to0$ e il gradiente diretto allo strato di compressione può
attenuarsi. L'esperimento 03 registra attivazioni e gradienti per monitorare questo
comportamento (\cref{ssec:hook}).

I pesi variazionali hanno forma $\shape{\nqlayers, \nq, 3}$. La classe
\classe{TorchLayer} li espone come parametri PyTorch, aggiornati da AdamW e
inclusi nei conteggi del \cref{cap:risultati}.

\subsection{Il ponte fra PyTorch e PennyLane}
\label{sec:integrazione}

Il QNode usa l'interfaccia PyTorch e la differenziazione \code{backprop};
\classe{TorchLayer} lo incapsula in un \code{nn.Module}. Il grafo di
autograd comprende così i pesi variazionali e gli strati classici a monte.

I dispositivi \code{default.qubit} lavorano su CPU. \classe{PinnedTorchLayer}
mantiene i pesi sul dispositivo del simulatore e riporta l'output sul dispositivo
del chiamante. \classe{QuantumLinear} appiattisce le dimensioni che precedono
l'ultima e le ripristina dopo il circuito, accettando sia
$\shape{\dbatch, \dseq, \nq}$ sia $\shape{\dbatch, \npatch, \nq}$.

\begin{limite}
La simulazione e la gestione dei batch producono l'overhead misurato nel
\cref{cap:risultati}; non è stato svolto un profiling che separi i due contributi.
I tempi riguardano il simulatore e non stimano l'esecuzione su hardware
(\cref{sec:simulazione}).
\end{limite}

\section{L'adapter e la sua ambiguità}
\label{sec:adapter}

La libreria definisce \classe{QuantumLinearWithAdapter}, importata
dall'esperimento 01 come \classe{QuantumLayerAdapter}. La classe antepone al
circuito una proiezione da \code{n\_in} a \code{n\_qubits} e riusa internamente
\classe{QuantumLinear}.
Nel preset definitivo dell'esperimento~01 la dimensione di embedding è
$\demb = 8$, mentre il circuito usa $\nq = 4$ qubit; l'adapter rende compatibili
le due dimensioni. Negli
esperimenti 02 e 03 non serve, perché $\dvit = \nq$ è imposto per costruzione.

\begin{limite}
L'adapter è uno strato classico addestrabile posto immediatamente prima del
circuito. Qualunque effetto misurato nell'esperimento~01 è perciò l'effetto della
\emph{composizione} adapter${}+{}$VQC, e le due componenti non sono separabili
sulla base dei dati raccolti.
\end{limite}

Negli esperimenti 02 e 03 la proiezione classica è invece appaiata fra le
varianti. Nel 03 i tre modelli condividono la stessa
\code{nn.Linear(147, 10)} e differiscono nelle trasformazioni successive.

\section{Configurazioni e ciclo di addestramento}
\label{sec:base-config}

\classe{BaseViTConfig} raccoglie i parametri comuni agli esperimenti 02 e 03.
\classe{QViTConfig} usa \code{embed\_dim=8}; \classe{AblationConfig} lo porta a
10 e aggiunge le impostazioni dell'ablazione e della diagnostica.
\classe{GPTConfig} resta separata. I valori completi sono
nell'appendice~\ref{app:configurazioni}.

\code{n\_qubits} è una proprietà derivata da \code{embed\_dim}; in
\classe{GPTConfig} deriva invece da $\demb/\nhead$. \code{to\_dict()} include
questi valori nel \file{config.json}. Le verifiche sulle configurazioni sono
elencate nell'appendice~\ref{app:riproduzione}.

\subsection{\texorpdfstring{\classe{BaseTrainer}}{BaseTrainer}}
\label{sec:base-trainer}

\classe{BaseTrainer} incapsula il ciclo comune agli esperimenti 02 e 03:
ottimizzazione, valutazione, checkpointing, logging su TensorBoard, diagnostica e
salvataggio dei risultati. L'ottimizzatore è \code{AdamW} con
\code{learning\_rate} e \code{weight\_decay} presi dalla configurazione, lo
scheduler è \code{CosineAnnealingLR} con \code{T\_max = max\_epochs}, e il
\emph{gradient clipping} usa soglia $1{,}0$. La validation segue la cadenza della
configurazione e include sempre la prima e l'ultima epoca. Il checkpoint con
accuracy di validation più alta è salvato in \file{best\_model.pth}, mentre lo
stato dell'ultima epoca va in \file{final\_model.pth}.

\begin{limite}
Le metriche finali, inclusa l'accuracy di training usata nel gap di
generalizzazione, sono calcolate sul checkpoint selezionato sulla validation e
non sullo stato dell'ultima epoca.
\end{limite}

\subsection{Hook diagnostici}
\label{ssec:hook}

Nell'esperimento~03 gli \emph{hook} osservano
\code{"patch\_embed.compression"}, lo strato lineare appaiato fra le tre
varianti.

L'hook all'indietro, registrato con \code{register\_full\_backward\_hook},
memorizza a ogni passo la norma $L^2$ del gradiente in uscita dal modulo di
compressione. La misura può segnalare un'attenuazione lungo il percorso
successivo, ma non ne identifica da sola la causa. L'hook in avanti, registrato
con \code{register\_forward\_hook}, memorizza media,
deviazione standard, massimo assoluto e \code{saturation\_frac} delle attivazioni
prodotte. La frazione di saturazione è la frazione di valori con
\begin{equation}
  \abs{x} \;>\; \arctanh(0{,}95) \;\approx\; 1{,}83 ,
  \label{eq:saturazione}
\end{equation}
cioè quelli per cui $\abs{\tanh(x)} > 0{,}95$. Il modulo osservato produce
l'attivazione prima della $\tanh$; $0{,}95$ è la soglia diagnostica adottata.

\section{Metriche, statistica e riproducibilità}
\label{sec:analisi-statistica}

La funzione \code{compute\_classification\_metrics(labels, preds, probs)}
restituisce accuracy, macro-AUC \emph{one-vs-rest} e macro-F1, gestendo a parte il
caso binario attraverso la probabilità della classe positiva; se una partizione
contiene una sola classe l'AUC non è calcolabile e la funzione usa $0{,}0$ come
fallback tecnico. Quel valore non rappresenta una misura interpretabile e non si
verifica nelle partizioni definitive. La media \emph{macro} è preferita alla
\emph{micro} per lo sbilanciamento dei dataset MedMNIST, perché assegna lo stesso
peso a ogni classe.
\code{generalization\_gap} restituisce la differenza con segno
$\acctr - \accte$, che è la metrica primaria dell'esperimento~03.

Il confronto appaiato è affidato a \code{paired\_comparison()}, che calcola
differenze per seed, media e deviazione standard campionaria, intervallo $t$ al
\SI{95}{\percent}, paired $t$-test, intervallo bootstrap percentile su
\num{10000} ricampionamenti con generatore pseudocasuale fissato, test di Wilcoxon
dei ranghi con segno e dimensione dell'effetto $\dz = \dmean/\sd$. Con $n<2$ i
test non sono calcolati; con $\sd=0$ l'intervallo coincide con $\dmean$. Se tutte
le differenze sono nulle, il campo del Wilcoxon rimane \code{None}. La
giustificazione metodologica è nella
\cref{sec:statistica}.

\subsection{Provenienza e schema dei risultati}
\label{sec:riproducibilita}

\code{set\_seed(seed)} imposta i generatori di \code{random}, \code{numpy} e
\code{torch} e, quando una GPU è disponibile, anche quelli CUDA, attivando
\code{cudnn.deterministic}. Il codice non forza però algoritmi deterministici per
tutte le operazioni PyTorch; non sono state eseguite repliche dedicate per
misurare eventuali differenze numeriche.

\code{collect\_run\_metadata()} registra stato Git, comando, timestamp, versioni
delle librerie e hardware; il flag \emph{dirty} segnala modifiche non committate.
I JSON riportano inoltre la versione dello schema e rifiutano valori NaN. La
versione permette di riconoscere le run storiche prive di campi aggiunti in fasi
successive; il loro uso è discusso nella \cref{sec:minacce}.

La suite in \file{tests/} copre configurazioni, metriche, statistica, appaiamento
dei modelli, compatibilità fra simulatori e ripresa dell'addestramento. Non è una
verifica esaustiva dei modelli o delle run complete, per le quali il protocollo
prevede anche uno \emph{smoke test} a poche epoche
(appendice~\ref{app:riproduzione}).
